\chapter{Sensitivity to the steam-atmosphere model}
\label{chap:steam_sensitivity}

The baseline forecast of \Cref{chap:results} adopts the AGNI-updated
steam mass--radius grid with the Dorn \& Lichtenberg (2021) water-retention
correction as the fiducial physical prescription.  This chapter tests how
strongly the \HZIED detection depends on that choice by re-running the
same P4 and P5 analyses under each of the alternative steam-atmosphere
models introduced in \Cref{sec:runaway_models}.  Every sweep uses the
observable cuts of the AGNI baseline; only the mass--radius prescription
varies.  Figures in this chapter are produced by
\texttt{Src/figures/ch08\_steam\_sensitivity.py}.


% §HZIED Baseline: Turbet Steam Model moved to condensed §3.6
% (condensed_03_results.tex); labels live in 07_results.tex and condensed wrapper
\phantomsection\label{ssec:relaxed_depth}


% =======================================================================
\section{Dependence on the Steam Atmosphere Model}
\label{sec:steam_model}

The amplitude and shape of the predicted radius discontinuity depend
directly on the steam-atmosphere model used to inflate hot-side planets.
Five prescriptions beyond the Turbet baseline are compared here, all run
at $\texttt{depth\_min\_ppm} = 60$, $x_{\mathrm{H_2O}} = 0.005$,
$f_{\mathrm{rgh}} = 0.8$, and \texttt{age\_mode = `sampled'}.
\Cref{tab:steam_model_summary} gives an overview of all results.

\begin{table}[H]
    \centering
    \caption[Steam model comparison: planet yield and Bayes factor]{%
        Planet yield and Bayes factor for each steam-atmosphere prescription
        on the P4 and P5 samples (depth $\geq 60\,\mathrm{ppm}$, $x_{\mathrm{H_2O}}=0.005$,
        $f_{\mathrm{rgh}}=0.8$).  $N$ is the number of planets entering the
        \HZIED analysis; cold/infl.\ are planets on the cold and inflated
        hot side of $S_\mathrm{thresh}$.}
    \label{tab:steam_model_summary}
    \small
    \setlength{\tabcolsep}{5pt}
    \begin{tabular}{llrrrrr}
        \toprule
        Model & Sample & $M$ range ($M_\oplus$) & $N$ & cold & infl. & $\ln\Delta Z$ \\
        \midrule
        Turbet     & P4 & [0.1,\,2.05] &  80 & 27 & 44 &  24.5 \\
                   & P5 &              &  82 & 16 & 54 &  44.3 \\
        \midrule
        AGNI       & P4 & [0.1,\,10]   & 135 & 41 & 78 &  17.9 \\
                   & P5 &              & 172 & 26 &118 &  75.5 \\
        \midrule
        AGNI updated & P4 & [0.5,\,20] & 139 & 38 & 80 &  44.9 \\
                   & P5 &              & 266 & 33 &186 &  15.6 \\
        \midrule
        Aguichine  & P4 & [0.2,\,20]   & 159 & 46 & 90 &  49.2 \\
                   & P5 &              & 270 & 37 &185 &  31.7 \\
        \midrule
        Pierrehumbert & P4 & [1.5,\,9] &  64 & 24 & 35 & 149.9 \\
        ($S_\mathrm{thresh}=435\,\mathrm{W\,m^{-2}}$) & P5 & & 127 & 16 & 88 &  35.9 \\
        \midrule
        Hycean mixed & P4 & [0.5,\,20] & 139 & 38 & 75 &  75.0 \\
        ($f_\mathrm{hycean}=0.5$) & P5 &  & 277 & 33 &207 &  93.9 \\
        \bottomrule
    \end{tabular}
\end{table}


\subsection{Turbet steam table (reference)}
\label{ssec:turbet}

The \textcite{Turbet2020} model is a 1D grey-atmosphere inverse
radiative-convective model applied to a rocky baseline from
\textcite{Zeng2016}.  The steam-atmosphere radius increment is tabulated
as a function of planet mass ($M \in [0.1,\,2.05]\,\Mearth$) and water
mass fraction ($x_{\mathrm{H_2O}} \in [0,\,0.05]$).  This is the
Schlecker+2024 fiducial and serves as the reference for all comparisons
in this section.  Results at $60\,\mathrm{ppm}$ are reported in
\Cref{ssec:relaxed_depth}: P4 $N = 80$, $\ln\Delta Z = 24.5$;
P5 $N = 82$, $\ln\Delta Z = 44.3$.


\subsection{AGNI grid}
\label{ssec:agni}

The AGNI radiative-transfer grid (grid\_022, \textcite{BoerNichollsLichtenberg2025})
tabulates the photospheric radius of a steam-bearing planet on a 2D grid
of planet mass and instellation.  For $M < 1\,\Mearth$ the grid falls
back to the \textcite{Turbet2020} table.  The combined observable range
is $M \in [0.1,\,10]\,\Mearth$, extending to five times the Turbet
ceiling and capturing the super-Earth population.

The wider mass range substantially increases the planet yield: P4 gains
$+69\%$ planets ($N = 135$) and P5 gains $+110\%$ ($N = 172$) relative
to the Turbet reference.  Despite more planets, the P4 Bayes factor drops
to $\ln\Delta Z = 17.9$: the additional super-Earth-mass planets at
$M \in [2,\,10]\,\Mearth$ experience smaller fractional radius inflation
from the AGNI model at the fiducial $x_{\mathrm{H_2O}} = 0.005$, diluting
the step signal.  P5 benefits strongly ($\ln\Delta Z = 75.5$) because its
larger stellar radii mean that the deeper transits of heavier planets
bring more objects above the $60\,\mathrm{ppm}$ cut.

\begin{figure}[H]
    \centering
    \includegraphics[width=\textwidth]{plato_agni_hzied}
    \caption[HZIED with AGNI steam model]{%
        Detected planet sample for P4 (left) and P5 (right) using the
        AGNI grid \parencite{BoerNichollsLichtenberg2025},
        $M \in [0.1,\,10]\,\Mearth$, depth $\geq 60\,\mathrm{ppm}$.
        P4: $N=135$, $\ln\Delta Z = 17.9$.
        P5: $N=172$, $\ln\Delta Z = 75.5$.}
    \label{fig:plato_agni_hzied}
\end{figure}


\subsection{Extended AGNI + PALEOS (\texttt{agni\_updated})}
\label{ssec:agni_updated}

The \texttt{agni\_updated} prescription \parencite{BoerNichollsLichtenberg2025}
replaces the 2D grid with a 1D AGNI+PALEOS steam table computed via the
ratio method: the inflated radius is obtained by scaling the PALEOS rocky
baseline by the AGNI steam-to-dry radius ratio, ensuring consistency
between the interior and atmosphere models.  The \textcite{DornLichtenberg2021}
water-retention correction is applied on top.  The table covers
$M \in [0.5,\,20]\,\Mearth$; below $0.5\,\Mearth$ the stratosphere
temperature is prescribed rather than solved, so lower-mass planets are
excluded.

Opening the mass range to $20\,\Mearth$ doubles the P5 yield to $N = 266$
($+225\%$ vs.\ Turbet) and nearly doubles P4 to $N = 139$.  However, the
P5 Bayes factor falls to $\ln\Delta Z = 15.6$: at the fiducial water
fraction, the large-mass tail of the \texttt{agni\_updated} grid predicts
modest inflation, and the many additional mini-Neptune-mass planets add
noise without sharpening the step.  P4 is more robust ($\ln\Delta Z = 44.9$)
because M~dwarfs predominantly host lower-mass planets where the inflation
amplitude is larger relative to the scatter.

\begin{figure}[H]
    \centering
    \includegraphics[width=\textwidth]{plato_agni_updated_hzied}
    \caption[HZIED with AGNI updated steam model]{%
        Detected planet sample for P4 (left) and P5 (right) using the
        \texttt{agni\_updated} prescription \parencite{BoerNichollsLichtenberg2025},
        $M \in [0.5,\,20]\,\Mearth$, depth $\geq 60\,\mathrm{ppm}$.
        P4: $N=139$, $\ln\Delta Z = 44.9$.
        P5: $N=266$, $\ln\Delta Z = 15.6$.}
    \label{fig:plato_agni_updated_hzied}
\end{figure}


\subsection{Aguichine irradiated ocean planets}
\label{ssec:aguichine}

\textcite{Aguichine2021} compute internal-structure models for irradiated
ocean planets consisting of a rocky core, a supercritical water layer,
and a steam atmosphere, tabulated as a function of core mass fraction,
water mass fraction (WMF), irradiation temperature $T_\mathrm{eq}$, and
--- unlike every steam prescription used earlier in this chapter ---
planet \emph{age} in Gyr.  The model is physically distinct from the
AGNI family: rather than treating the steam layer as a thin atmospheric
increment on a rocky baseline, it self-consistently solves for the bulk
radius of a water-dominated interior that contracts with time as
gravitational-energy release and radiogenic heat diminish.  The
applicable mass range is $M \in [0.2,\,20]\,\Mearth$, and the tabulated
age range spans $0.1$--$20\,\mathrm{Gyr}$.

At the fixed benchmark configuration used up to this point in the
chapter --- \texttt{aguichine} steam model, $x_{\mathrm{H_2O}}=0.005$,
$f_\mathrm{rgh}=0.8$, $S_\mathrm{thresh}=280\,\mathrm{W\,m^{-2}}$,
depth $\geq 60\,\mathrm{ppm}$, $M \in [0.2,\,20]\,\Mearth$, and a
strict FLAME age cut --- the planet yields are comparable to
\texttt{agni\_updated} (P4 $N \approx 160$--$170$, P5 $N \approx 270$).
The Bayes factors are also comparable
($\ln\Delta Z \sim 20$--$50$), with the exact number depending
sensitively on the age assignment scheme (see below) and on the
noise-realisation seed
(the seed-to-seed jitter is a factor of ${\sim}2$ in $\ln\Delta Z$,
as quantified in \Cref{ssec:wrr_sweep} for the analogous
\texttt{agni\_updated} sweep).

\paragraph{Context: what has been kept fixed so far.}
Every subsection of \Cref{sec:steam_model} up to this point has used
\texttt{age\_mode='sampled'} --- FLAME ages where Gaia provides them,
bootstrap-filled from the empirical FLAME distribution for the
$\sim 21\,\%$ of P5 stars and $\sim 76\,\%$ of P4 stars without a
FLAME age (see \Cref{sec:plato_samples} for the FLAME coverage).
That choice preserved the full stellar pool while keeping the age
distribution close to what Gaia actually measures.  The Aguichine
prescription is however the first steam model whose radius depends
on age; the choice of age-assignment scheme therefore feeds through
to the inflated radii and, in turn, to the recovered $\ln\Delta Z$
and $\hat{S}_\mathrm{thresh}$.  The following analysis quantifies
this sensitivity.

\paragraph{Age-assignment schemes compared.}
Four schemes are compared here:
\begin{itemize}
    \item \texttt{strict} --- keep only stars with a Gaia FLAME age
    (Creevey et al.\ 2023); drop the rest.  Physically clean but
    strongly reduces $N$, especially for P4 where FLAME coverage is
    ${\sim}24\,\%$.
    \item \texttt{sampled} --- FLAME where available; bootstrap-fill
    missing ages from the empirical FLAME distribution of the same
    sample.  Preserves $N$; assumes the unmeasured stars share the
    same age distribution as the measured ones.  This is the default
    used elsewhere in the chapter.
    \item \texttt{flame\_flat} --- FLAME where available; fill missing
    ages from a uniform $[0,\,10]\,\mathrm{Gyr}$ prior (mode added
    for this analysis).  Combines observed FLAME data with a
    maximally agnostic prior for the unmeasured tail --- important
    if one worries that FLAME's detection sensitivity biases its
    reported age distribution.
    \item \texttt{flat} --- uniform $[0,\,10]\,\mathrm{Gyr}$ for every
    star; FLAME ignored.  Reproduces the pre-FLAME pipeline
    behaviour and serves as an upper bound on how much age
    information can matter for the HZIED signal.
\end{itemize}

\begin{table}[H]
    \centering
    \caption[Fixed parameters of the Aguichine age analysis]{%
        Configuration held fixed while the age-assignment mode is varied.
        Only \texttt{age\_mode} changes across the sweep; every other
        parameter matches the baseline used in
        \Cref{sec:steam_model,sec:composition_dependence}.}
    \label{tab:aguichine_age_params}
    \small
    \begin{tabular}{lll}
        \toprule
        Parameter & Value & Note \\
        \midrule
        \multicolumn{3}{l}{\textit{Injection / model}} \\
        \texttt{radius\_inflation} & \texttt{aguichine} & Aguichine+2025 SWE table \\
        \texttt{use\_dl21}         & \texttt{True}      & Dorn \& Lichtenberg (2021) water incorporation \\
        $x_\mathrm{H_2O}$          & $0.005$            & Schlecker+2024 fiducial \\
        $S_\mathrm{thresh}$        & $280\,\mathrm{W\,m^{-2}}$ & rocky HZ inner edge \\
        $f_\mathrm{rgh}$           & $0.8$              & runaway greenhouse fraction \\
        \multicolumn{3}{l}{\textit{Observable cuts}} \\
        $M$                        & $[0.2,\,20]\,\Mearth$  & Aguichine table range \\
        depth                      & $\geq 60\,\mathrm{ppm}$ & relaxed cut (\Cref{ssec:relaxed_depth}) \\
        \multicolumn{3}{l}{\textit{Inference}} \\
        window $w$                 & $25$               & rolling SMA \\
        $\sigma_R / R$             & $4\%$ (P4), $3\%$ (P5) & PLATO Red Book \\
        $\sigma_S / S$             & $5\%$              & Schlecker+2024 \\
        \texttt{sma\_x\_mode}      & \texttt{standard}  & mean $S$ per window \\
        \texttt{sigma\_sma\_mode}  & \texttt{empirical} & $\mathrm{std}(R_\mathrm{obs})/\sqrt{w}$ \\
        \texttt{nlive}             & $100$              & dynesty live points \\
        seed base                  & $42+200\,i$ ($i=0..3$) & unique per age mode; 10 trials each (seeds $+13k$, $k=0..9$) \\
        \bottomrule
    \end{tabular}
\end{table}

\paragraph{Results.}
\Cref{fig:plato_aguichine_age_lnk} shows the recovered $\ln\Delta Z$
(left) and posterior-median $\hat{S}_\mathrm{thresh}$ (right) across
the four age modes for P4 and P5.  Each marker is the median across
$10$ independent noise realisations; the shaded band is the
$16$--$84\,\%$ range, capturing both measurement-noise scatter and
dynesty variance at $n_\mathrm{live}=100$ (same construction as
\Cref{ssec:wrr_sweep}).  The multi-seed medians (bracketed by the
per-seed min/max across the 10 draws) are collected in
\Cref{tab:aguichine_age_results}.

\begin{table}[H]
    \centering
    \caption[Multi-seed medians of the aguichine age sweep]{%
        Multi-seed medians of the aguichine age-mode sweep.  $\ln\Delta Z$
        columns quote the median across $10$ noise realisations with the
        per-seed min/max bracket; $\hat{S}_\mathrm{th}$ columns are the
        median of the posterior-median values across the same $10$ draws
        (W\,m$^{-2}$).  The injected threshold is $280\,\mathrm{W\,m^{-2}}$.}
    \label{tab:aguichine_age_results}
    \small
    \begin{tabular}{lccccccc}
        \toprule
        \texttt{age\_mode} & $N_\mathrm{P4}$ & P4 $\ln\Delta Z$ (med [min,max])
        & P4 $\hat{S}_\mathrm{th}$
        & $N_\mathrm{P5}$ & P5 $\ln\Delta Z$ (med [min,max])
        & P5 $\hat{S}_\mathrm{th}$ \\
        \midrule
        \texttt{strict}     & $31$  & $\phantom{-}0.1\;[-0.3,\,0.6]$   & $584$ & $173$ & $23.6\;[19.4,\,34.1]$  & $348$ \\
        \texttt{sampled}    & $159$ & $55.0\;[34.6,\,79.4]$            & $208$ & $270$ & $26.4\;[22.1,\,36.6]$  & $256$ \\
        \texttt{flame\_flat}& $148$ & $-1.3\;[-1.7,\,-0.3]$            & $256$ & $238$ & $42.2\;[35.9,\,56.2]$  & $359$ \\
        \texttt{flat}       & $159$ & $62.4\;[48.7,\,71.9]$            & $205$ & $238$ & $38.6\;[29.3,\,52.6]$  & $385$ \\
        \bottomrule
    \end{tabular}
\end{table}

Unlike the wrr and Hycean sweeps, this one exposes a real
age-driven sensitivity in the aguichine model, and the four modes
respond in genuinely different ways.

The \texttt{strict} mode collapses on P4: only about a quarter of
M-dwarf hosts carry a FLAME age (\Cref{sec:plato_samples}), so
requiring one drops the sample to $N=31$ and the step-function
inference has no signal to find --- $\ln\Delta Z\approx 0$ and the
recovered $\hat{S}_\mathrm{thresh}=584\,\mathrm{W\,m^{-2}}$ is
prior-dominated.  On P5 the same cut is far less damaging (FLAME
covers roughly $79\,\%$ of stars), and \texttt{strict} still yields
$\ln\Delta Z\approx 24$ on $N=173$ planets.

The \texttt{sampled} mode recovers a strong signal on P4
($\ln\Delta Z\approx 55$) because bootstrapping from the observed
FLAME distribution preserves the narrow, mostly-old age pool that
the M-dwarf hosts already sit in.  Old planets on the hot side
contract more in the aguichine model, so the step at
$S_\mathrm{thresh}$ stays coherent and the sampler picks it up.

The \texttt{flame\_flat} mode is the most instructive of the four
because it actually \emph{disfavours} H\textsubscript{1} on P4,
returning $\ln\Delta Z = -1.3$.  The mechanism is a bimodal age
distribution: the $24\,\%$ of stars with a real FLAME age sit in a
narrow, old band, while the $76\,\%$ filled from the uniform prior
spread out over $[0,\,10]\,\mathrm{Gyr}$.  On the hot side of
$S_\mathrm{thresh}$ that bimodality broadens the observed radius
distribution enough to wash out the step, and the posterior tips
slightly toward the flat model.  P5 shows the mirror-image effect:
FLAME covers most of the sample there, so the uniform fill only
touches $\approx 21\,\%$ of the stars, and its younger tail
\emph{sharpens} the step by adding a handful of strongly-inflated
hot-side planets.

The \texttt{flat} mode --- uniform ages everywhere --- gives the
strongest P4 signal ($\ln\Delta Z\approx 62$) because a uniform
distribution keeps the sample-averaged age around $5\,\mathrm{Gyr}$,
slightly younger than the FLAME pool, and produces a larger typical
inflation ratio without any of the bimodality that hurts
\texttt{flame\_flat}.  On P5 it is statistically indistinguishable
from \texttt{flame\_flat}.

The recovered $\hat{S}_\mathrm{thresh}$ pattern (right panel of
\Cref{fig:plato_aguichine_age_lnk}) also carries information: on P4
the two modes with strong evidence (\texttt{sampled} and \texttt{flat})
recover $\hat{S}\approx 205$--$208\,\mathrm{W\,m^{-2}}$, undershooting
the injected $280$ by ${\sim}70\,\mathrm{W\,m^{-2}}$ --- the same
sub-fiducial pull seen in the wrr sweep at low $x_{\mathrm{H_2O}}$
(\Cref{ssec:wrr_sweep}).  On P5 the recovery straddles the injected
value, from $256$ (\texttt{sampled}) to $385$ (\texttt{flat}).  The
strict-P4 outlier ($\hat{S}=584$) and flame\_flat-P4
($\hat{S}=256$ but $\ln\Delta Z\approx 0$) both reflect the H\textsubscript{0}-preferred
regime and should not be interpreted as physical thresholds.

The practical take-away for the wider chapter is that the
$\texttt{sampled}$ mode adopted elsewhere in the chapter (see
\Cref{sec:steam_model,sec:composition_dependence}) is a reasonable
default for the AGNI-family models --- where the radius does not depend
on age --- but is not the only defensible choice for aguichine and
similar age-sensitive prescriptions.  A future analysis of an
aguichine-dominated population would benefit from either
$\texttt{strict}$ (physically clean but sample-limited on P4) or a
population-level marginalisation over the age uncertainty.

\Cref{fig:plato_aguichine_age_panels} shows the detected planet sample
for one representative trial per (sample, age mode), with the
per-planet grey inflation lines from the rocky baseline to the
Aguichine radius.  The visual differences between the modes are
subtle: the underlying stellar catalogue is the same; only the
assigned ages differ, and the Aguichine radius varies only
${\sim}0.05\,\Rearth$ across the tabulated 0.1--20 Gyr range for
typical intermediate-mass planets.

\begin{figure}[H]
    \centering
    \includegraphics[width=\textwidth]{plato_aguichine_age_lnk}
    \caption[Aguichine age-mode sweep: $\ln\Delta Z$ and $\hat{S}_\mathrm{thresh}$]{%
        Left: Bayes factor $\ln\Delta Z$ as a function of the age-assignment
        mode for P4 (blue) and P5 (red) with the \texttt{aguichine} steam
        prescription.  Right: the corresponding posterior-median of the
        recovered runaway-GH threshold $\hat{S}_\mathrm{thresh}$; the
        dotted horizontal line is the injected value
        $280\,\mathrm{W\,m^{-2}}$.  In both panels each marker is the
        median of $10$ independent noise realisations and the shaded
        band spans the $16$--$84\,\%$ range of those draws.  All non-age
        pipeline parameters are as listed in
        \Cref{tab:aguichine_age_params}.}
    \label{fig:plato_aguichine_age_lnk}
\end{figure}

\begin{figure}[H]
    \centering
    \includegraphics[width=\textwidth]{plato_aguichine_age_panels}
    \caption[Aguichine age-mode sweep: detected planet samples]{%
        Detected planet sample for P4 (top row) and P5 (bottom row) at
        each of the four age-assignment modes (columns; representative
        trial per mode).  Vertical grey lines connect each planet's
        rocky baseline radius to its Aguichine-inflated radius,
        illustrating the small age-dependent contraction envelope
        (${\sim}0.05\,\Rearth$ over $0.1$--$20\,\mathrm{Gyr}$ for
        intermediate-mass planets).  Titles quote $N$, the cold /
        inflated planet counts, and the recovered $\ln\Delta Z$ for
        the shown trial.}
    \label{fig:plato_aguichine_age_panels}
\end{figure}

% Figures produced by:
%   Src/figures/plato_aguichine_age_analysis.py
%     → Figures/generated/plato_aguichine_age_lnk.pdf
%     → Figures/generated/plato_aguichine_age_panels.pdf
%
% The original single-seed representative scatter figure (Fig 6.6 above)
% remains as `plato_aguichine_hzied` and is produced by
%   Src/figures/plato_steam_models_hzied.py.


\subsection{Pierrehumbert Hycean model}
\label{ssec:pierrehumbert}

\textcite{Pierrehumbert2022} describes runaway greenhouse conditions for
ocean planets with a hydrogen-dominated atmosphere overlying a
condensed-water or supercritical-water interior — so-called Hycean
worlds.  The runaway greenhouse threshold for this scenario lies at a
substantially higher instellation than for rocky planets:
$S_\mathrm{thresh} = 435\,\mathrm{W\,m^{-2}}$ \parencite{Innes2023},
compared to $280\,\mathrm{W\,m^{-2}}$ for the rocky prescriptions.
The applicable mass range is $M \in [1.5,\,9]\,\Mearth$.

The higher $S_\mathrm{thresh}$ and narrower mass range reduce the sample
sizes ($N = 64$ for P4, $N = 127$ for P5), but the \HZIED signal is
exceptionally strong: P4 yields $\ln\Delta Z = 149.9$, the highest of all
models tested.  The origin is the large predicted radius inflation for
Hycean worlds — a planet crossing the $435\,\mathrm{W\,m^{-2}}$ threshold
swells from a dense water-rich interior to a fully steam-dominated
configuration, producing a radius step far larger than in the rocky
models.  P5 is also strongly detected ($\ln\Delta Z = 35.9$).

\begin{figure}[H]
    \centering
    \includegraphics[width=\textwidth]{plato_hycean_pierrehumbert_hzied}
    \caption[HZIED with Pierrehumbert Hycean model]{%
        Detected planet sample for P4 (left) and P5 (right) using the
        \textcite{Pierrehumbert2022} Hycean prescription,
        $M \in [1.5,\,9]\,\Mearth$, $S_\mathrm{thresh} = 435\,\mathrm{W\,m^{-2}}$,
        depth $\geq 60\,\mathrm{ppm}$.
        P4: $N=64$, $\ln\Delta Z = 149.9$.
        P5: $N=127$, $\ln\Delta Z = 35.9$.}
    \label{fig:plato_hycean_pierrehumbert_hzied}
\end{figure}


\subsection{Hybrid rocky--Hycean model}
\label{ssec:hycean_mixed}

The hybrid model introduces two free parameters that together encode the
unknown Hycean prevalence in the planet population.

\paragraph{Free parameters.}
\begin{itemize}
    \item $f_\mathrm{hycean} \in [0,1]$: the fraction of planets in the
    Hycean mass range $M \in [1.5,\,9]\,\Mearth$ that are treated as
    Hycean worlds.  The complementary fraction $1-f_\mathrm{hycean}$
    follows the \texttt{agni\_updated} rocky prescription.
    $f_\mathrm{hycean} = 0$ recovers the pure \texttt{agni\_updated} result;
    $f_\mathrm{hycean} = 1$ is the maximally Hycean scenario.
    \item $S_\mathrm{thresh,\,Hycean}$: the runaway greenhouse instellation
    threshold for Hycean planets.  Following \textcite{Innes2023}, this
    is set to $435\,\mathrm{W\,m^{-2}}$, substantially higher than the
    rocky threshold of $280\,\mathrm{W\,m^{-2}}$.  Planets with
    $S_\mathrm{abs} \in [280,\,435]\,\mathrm{W\,m^{-2}}$ are therefore
    classified differently depending on whether they are rocky or Hycean:
    hot and inflated under the rocky model, but cool and uninflated under
    the Hycean model.
\end{itemize}

The total observable mass range is $M \in [0.5,\,20]\,\Mearth$, with the
rocky \texttt{agni\_updated} model applying to all planets and the Hycean
override applying to the randomly selected $f_\mathrm{hycean}$ fraction
of $M \in [1.5,\,9]\,\Mearth$ planets.

\paragraph{Dependence on $f_\mathrm{hycean}$.}
\Cref{fig:plato_hycean_fraction_lnk} shows how the Bayes factor
$\ln\Delta Z$ (left panel) and the recovered runaway-GH threshold
$\hat{S}_\mathrm{thresh}$ (right panel) vary with $f_\mathrm{hycean}$
for P4 and P5.  Each marker is the median across $10$ independent
noise realisations and the shaded band gives the 16--84\,\% range,
capturing the combined jitter from measurement-noise draws and dynesty
variance at $n_\mathrm{live}=100$ (see \Cref{ssec:wrr_sweep} for the
same construction on the water-content sweep).

For both samples the median evidence rises monotonically with
$f_\mathrm{hycean}$.  On P4 the growth is modest (from
$\ln\Delta Z \approx 49$ at $f_\mathrm{hycean}=0$ to $\approx 95$ at
$f_\mathrm{hycean}=1$, a factor of $\approx 2$), because M-dwarf HZ
planets sit at much smaller absolute instellations than $435\,\mathrm{W\,m^{-2}}$
and only a modest number are reclassified when the Hycean threshold is
imposed.  On P5 the gain is dramatic (from $\approx 18$ to $\approx 216$,
a factor of $\approx 12$): a large fraction of FGK-star planets sits in
the reclassification window $S_\mathrm{abs}\in[280,\,435]\,\mathrm{W\,m^{-2}}$
and is shifted from the inflated to the uninflated side as
$f_\mathrm{hycean}$ grows, producing a sharper effective step at the
Hycean threshold that is easier to detect.

The recovered $\hat{S}_\mathrm{thresh}$ tells a complementary and
somewhat unexpected story.  At $f_\mathrm{hycean}=0$ both samples
recover values close to the injected rocky threshold
($\hat{S}_\mathrm{thresh}=246\,\mathrm{W\,m^{-2}}$ on P4,
$259\,\mathrm{W\,m^{-2}}$ on P5), with the same slight undershoot
seen in \Cref{ssec:wrr_sweep}.  As $f_\mathrm{hycean}$ increases the
posterior median migrates upward on both samples, but for
$f_\mathrm{hycean}\geq 0.5$ it settles \emph{well above} the injected
Hycean threshold of $435\,\mathrm{W\,m^{-2}}$: the P5 sample recovers
$\hat{S}_\mathrm{thresh}\approx 788\,\mathrm{W\,m^{-2}}$ at
$f_\mathrm{hycean}=0.5$ and $\approx 550$--$600\,\mathrm{W\,m^{-2}}$
at $f_\mathrm{hycean}\in\{0.75,\,1\}$; P4 similarly overshoots to
$\approx 590$--$880\,\mathrm{W\,m^{-2}}$.  This is not a failure of
recovery so much as a limit of the single-step hypothesis: for high
$f_\mathrm{hycean}$ the true population has \emph{two} effective steps
(the residual rocky one at $280\,\mathrm{W\,m^{-2}}$ and the Hycean one
at $435\,\mathrm{W\,m^{-2}}$) plus a systematically larger Hycean
baseline radius, and a single-parameter step function is forced to
compromise by placing itself where the local radius contrast against
the fitted baseline is maximised — typically hotter than either
injected threshold.  Practically this means that a decisive detection
($\ln\Delta Z \gg 5$) does \emph{not} imply an accurate
$\hat{S}_\mathrm{thresh}$ when the underlying composition mix is
non-trivial; a two-population Bayesian model would be required to
disentangle the two thresholds.

\Cref{fig:plato_hycean_fraction_panels} shows the detected planet
sample for $f_\mathrm{hycean} \in \{0,\,0.5,\,1\}$: the vertical grey
inflation lines lengthen and shift rightward as $f_\mathrm{hycean}$
increases, visually reflecting both the higher Hycean
$S_\mathrm{thresh}$ and the larger steam-atmosphere increment.

\begin{figure}[H]
    \centering
    \includegraphics[width=\textwidth]{plato_hycean_fraction_lnk}
    \caption[$\ln\Delta Z$ and recovered $\hat{S}_\mathrm{thresh}$ vs Hycean fraction]{%
        Left: Bayes factor $\ln\Delta Z$ as a function of the Hycean
        fraction $f_\mathrm{hycean}$ for P4 (blue) and P5 (red).
        Right: the corresponding posterior-median of the recovered
        runaway-GH threshold $\hat{S}_\mathrm{thresh}$; the two
        dotted/dashed horizontal lines mark the rocky value
        $280\,\mathrm{W\,m^{-2}}$ and the Hycean value
        $435\,\mathrm{W\,m^{-2}}$ \parencite{Innes2023}.  In both panels
        each marker is the median of $10$ independent noise realisations
        and the shaded band spans the 16--84\,\% range of those draws.
        Depth $\geq 60\,\mathrm{ppm}$, $M \in [0.5,\,20]\,\Mearth$.}
    \label{fig:plato_hycean_fraction_lnk}
\end{figure}

\begin{figure}[H]
    \centering
    \includegraphics[width=\textwidth]{plato_hycean_fraction_panels}
    \caption[Detected planet sample for key Hycean fractions]{%
        Detected planet sample for P4 (top row) and P5 (bottom row) at
        $f_\mathrm{hycean} \in \{0,\,0.5,\,1\}$ (left to right).
        Vertical grey lines connect each inflated planet's pre-inflation
        radius to its observed radius, showing the amplitude of the
        steam-atmosphere increment per model.
        $f_\mathrm{hycean}=0$ is equivalent to the pure
        \texttt{agni\_updated} prescription.}
    \label{fig:plato_hycean_fraction_panels}
\end{figure}

% Figures produced by:
%   Src/figures/plato_steam_models_hzied.py  →  Figures/generated/plato_*_hzied.pdf
%   Src/figures/plato_hycean_fraction_sweep.py
%       → Figures/generated/plato_hycean_fraction_lnk.pdf
%       → Figures/generated/plato_hycean_fraction_panels.pdf


